% ============================================================
%  MEAM 4600 — AI for Science and Engineering
%  Project 3 Technical Report
%  Sea-Level-Rise Adjusted Foundation Flood Risk Scorer
% ============================================================
\documentclass[conference]{IEEEtran}

\usepackage[T1]{fontenc}
\usepackage[utf8]{inputenc}
\usepackage{amsmath,amssymb}
\usepackage{graphicx}
\usepackage{url}
\usepackage{hyperref}
\usepackage{booktabs}
\usepackage{xcolor}
\usepackage{soul}          % \hl for highlighting
\usepackage{caption}

% ------------------------------------------------------------------
%  \placeholder{} — highlights a result that must be filled in
%  from actual run output before submission
% ------------------------------------------------------------------
\newcommand{\placeholder}[1]{%
  \colorbox{yellow!60}{\textbf{[#1]}}%
}

\graphicspath{{../outputs/}}

\title{Sea-Level-Rise Adjusted Foundation Flood Risk Scorer:\\
A GIS and Machine Learning Pipeline for Miami-Dade Coastal Parcels}

\author{%
  \IEEEauthorblockN{\placeholder{Student Name}}\\
  \IEEEauthorblockA{Department of Mechanical Engineering and Applied Mechanics\\
    University of Pennsylvania\\
    MEAM 4600 — AI for Science and Engineering\\
    \placeholder{email@seas.upenn.edu}}
}

\begin{document}
\maketitle

% ============================================================
\begin{abstract}
% ============================================================
Chronic coastal flooding driven by sea-level rise (SLR) poses a
compounding threat to the structural integrity of Miami-Dade County
condominiums, a risk made acutely visible by the 2021 Champlain Towers
South collapse and Florida Senate Bill 4-D milestone inspection mandate.
This report presents a reproducible, end-to-end pipeline—termed the
Sea-Level-Rise Adjusted Foundation Flood Risk Scorer (SIRS)—that
integrates FEMA National Flood Hazard Layer polygons, Miami-Dade County
parcel geometries, USGS 3DEP 1/3~arc-second elevation data, and NOAA
2022 SLR scenarios to quantify parcel-level chronic inundation risk.
A stillwater physics model computes the first year each parcel crosses
the chronic inundation threshold ($\text{MHHW} = 0.87$~ft~NAVD88), and
logistic regression produces an interpretable 0--100 risk score.
Under the NOAA Intermediate SLR scenario,
\placeholder{X}\% of studied Brickell parcels are projected to experience
chronic inundation before 2050.
The pipeline is designed as decision-support tooling for structural
engineering firms such as M2E Consulting Engineers performing SIRS
and milestone inspections under Florida SB~4-D.
\end{abstract}


% ============================================================
\section{Introduction}
% ============================================================

On June 24, 2021, Champlain Towers South in Surfside, Florida,
collapsed, killing 98~people and prompting the Florida Legislature to
enact Senate Bill 4-D, which mandates milestone structural inspections
for all condominium buildings three stories or taller within 25~years
of construction~\cite{floridaSB4D}.
Although the collapse's immediate cause was decades of deferred maintenance
and structural deterioration, subsequent engineering reviews highlighted
the role of chronic groundwater infiltration and soil subsidence—
conditions increasingly exacerbated by rising sea levels and intensifying
king-tide events along Miami's Biscayne Bay waterfront.

Chronic coastal flooding—defined as inundation that recurs ten or more
times per year~\cite{moftakhari2017}—differs from the episodic hurricane
flooding that dominates the FEMA regulatory framework.
Repeated wetting and drying accelerates corrosion of steel reinforcement
in post-tensioned slabs, degrades masonry mortar joints, and promotes
biological growth that compromises bearing capacity.
As NOAA SLR projections for the Virginia Key gauge station (8723214)
show between \placeholder{0.4} and \placeholder{3.2}~ft of rise by
2075 under Low and High scenarios respectively~\cite{sweet2022}, the
frequency and duration of chronic flooding events will increase
non-linearly for low-lying coastal parcels.

The scientific question motivating this work is whether
\textit{parcel-level chronic inundation onset timelines can be estimated
reproducibly from entirely public data sources}.
The goal is not to advance the state of the art in flood modeling—
physics-based hydrodynamic models exist for that—but to demonstrate
a GIS--ML coupling that an engineering firm can deploy using freely
available federal datasets and open-source Python tooling.
The result is an interpretable risk score that can be embedded directly
in a milestone inspection report or 30-year reserve study.

The remainder of this paper is organized as follows.
Section~\ref{sec:related} surveys related work in SLR science, compound
flooding, and ML-based risk assessment.
Section~\ref{sec:data} describes the four public data sources.
Section~\ref{sec:methods} details the spatial assembly, physics model,
and logistic regression scorer.
Section~\ref{sec:results} presents the outputs.
Section~\ref{sec:discussion} discusses engineering implications and
limitations, and Section~\ref{sec:conclusion} concludes.


% ============================================================
\section{Related Work}
\label{sec:related}
% ============================================================

\subsection*{Sea-Level Rise Science}

The primary authoritative source for near-term SLR projections at U.S.
tide gauges is the NOAA 2022 Sea-Level Rise Technical Report~\cite{sweet2022},
which provides five scenarios—Low, Intermediate-Low, Intermediate,
Intermediate-High, and High—calibrated to local relative sea-level trends
observed at individual gauge stations.
At Virginia Key (NOAA Station 8723214), subsidence amplifies global
mean SLR, producing a range from approximately 0.7~ft (Low) to
3.2~ft (High) above the year-2000 baseline by 2075.
Wdowinski et al.~\cite{wdowinski2016} demonstrated that repetitive
flooding frequency in Miami Beach correlates directly with observed
sea-level trends and is already detectable in the existing tidal record,
providing observational grounding for scenario-based projections.

\subsection*{Compound and Chronic Flooding}

Moftakhari et al.~\cite{moftakhari2017} formalized the distinction between
episodic extreme-event flooding and \textit{chronic} flooding, which they
define as inundation exceeding the Mean Higher High Water (MHHW) datum
ten or more times per year.
They show that chronic flood frequency is growing faster than extreme-event
frequency because even modest SLR shifts the tidal distribution upward
relative to fixed coastal infrastructure.
This paper adopts their threshold criterion: a parcel is chronically
flooded in year $t$ when its ground elevation falls below
$\text{MHHW} + \text{SLR}(t)$.

\subsection*{Machine Learning in Flood Risk Assessment}

Hauer et al.~\cite{hauer2016} projected that 13.1 million U.S.\ residents
face chronic inundation by 2100 under unchecked SLR, using census-block
population data coupled to NOAA inundation contours.
Their work established the population-scale framing but did not attempt
parcel-level resolution, SLR uncertainty propagation, or interpretable
risk scoring.
To our knowledge, no existing published pipeline couples parcel-level
GIS data to physics-based chronic inundation thresholds under multiple
SLR scenarios and then distills the result into a logistic regression
score designed for structural engineering decision support—
an interpretability-first gap that this work addresses.


% ============================================================
\section{Data Sources}
\label{sec:data}
% ============================================================

\begin{enumerate}

\item \textbf{FEMA National Flood Hazard Layer (NFHL).}
  Obtained from the FEMA Flood Map Service Center
  (\url{https://msc.fema.gov}) for FEMA FIRM panel 12086C
  (Miami-Dade County).
  Two layers were used: \texttt{S\_FLD\_HAZ\_AR} (flood hazard area
  polygons carrying the FLD\_ZONE classification and STATIC\_BFE
  attribute) and \texttt{S\_BFE} (base flood elevation contour lines
  carrying the BFE\_LEN attribute).
  The NFHL is in the public domain with no license restriction.

\item \textbf{Miami-Dade County Parcels.}
  Miami-Dade County parcel polygons with Florida DOR land-use codes
  were obtained from the Miami-Dade Open Data portal
  (\url{https://gis-mdc.opendata.arcgis.com}).
  The study area was clipped to the Brickell / Biscayne Bay waterfront
  bounding box (lon $-80.200$ to $-80.185$, lat $25.755$ to $25.775$)
  and filtered to residential and condominium land-use codes
  (DOR codes 01, 04, 08, 11).
  Data are released under the Miami-Dade County Open Data License
  (CC~BY~4.0 compatible).

\item \textbf{USGS 3D Elevation Program (3DEP) DEM.}
  A 1/3~arc-second (${\sim}10$~m) digital elevation model referenced to
  NAVD88 was downloaded from The National Map
  (\url{https://apps.nationalmap.gov/downloader}) covering the study area.
  Elevation was sampled at parcel centroids using \texttt{rasterio.sample()}.
  The 3DEP products are in the public domain (USGS).

\item \textbf{NOAA CO-OPS Virginia Key (Station 8723214).}
  Two products were used from the NOAA Center for Operational
  Oceanographic Products and Services (CO-OPS) API
  (\url{https://tidesandcurrents.noaa.gov}):
  (a)~hourly verified water-level observations for calendar year 2023
  (NAVD88 datum, GMT, English units) used to construct the tidal
  exceedance frequency curve, and
  (b)~the NOAA 2022 Sea-Level Rise Technical Report scenario
  values~\cite{sweet2022} for the Low, Intermediate, and High scenarios,
  hard-coded from Table A.5 at five-year anchor points from 2020 to 2075.
  All CO-OPS data are in the public domain.

\end{enumerate}


% ============================================================
\section{Methodology}
\label{sec:methods}
% ============================================================

\subsection{Spatial Data Assembly}

All spatial layers were standardized to EPSG:6437
(NAD83(2011) / Florida East, metre units) to ensure that distance
calculations and area computations are in metric units without
introducing the distortions of geographic coordinate systems.
The FEMA flood zone polygons were spatially joined to parcel centroids
using a \texttt{within} predicate in \texttt{geopandas.sjoin}.
When a centroid fell within multiple zones (due to polygon overlaps
at zone boundaries), zones were prioritized in the order
VE $>$ AE $>$ AH $>$ AO $>$ X (shaded) $>$ X (unshaded),
retaining only the highest-risk classification.

Base flood elevation (BFE) was assigned in a three-step cascade:
(1)~the STATIC\_BFE attribute from the matched flood zone polygon,
(2)~the BFE\_LEN value from the nearest \texttt{S\_BFE} line within
500~m using \texttt{geopandas.sjoin\_nearest}, or
(3)~the median BFE of all parcels in the same flood zone.
Ground-surface elevation in feet NAVD88 was sampled at each parcel
centroid from the 3DEP DEM after reprojecting centroid coordinates to
match the DEM's native geographic CRS using a \texttt{pyproj.Transformer}.
Parcels with missing elevation after sampling were filled using the
median elevation of the five spatially nearest parcels with valid
elevation via a \texttt{scipy.spatial.KDTree} nearest-neighbor search.

\subsection{Physics Model}

The stillwater inundation depth at parcel $x$ in year $t$ is:
\begin{equation}
  d(x, t) \;=\; \text{BFE} + \text{SLR}(t) - \text{elev}(x),
  \label{eq:depth}
\end{equation}
where BFE and elev$(x)$ are in feet NAVD88 and SLR$(t)$ is the
cumulative sea-level rise above the year-2000 baseline interpolated
linearly between the NOAA anchor-year values.
Positive $d$ indicates inundation depth during a 100-year flood event;
negative $d$ indicates remaining freeboard above the flood level.

A parcel is defined as \textit{chronically inundated} in year $t$ when
its ground elevation falls below the combined MHHW and SLR surface:
\begin{equation}
  \text{elev}(x) < \text{MHHW} + \text{SLR}(t),
  \label{eq:chronic}
\end{equation}
where MHHW~$= 0.87$~ft~NAVD88 at Virginia Key Station 8723214,
as tabulated in the NOAA CO-OPS datum table~\cite{noaa_datums}.
The chronic inundation onset year $T_c(x)$ is defined as the first
year in the study period $[2025, 2075]$ for which
condition~\eqref{eq:chronic} is satisfied:
\begin{equation}
  T_c(x) \;=\; \min\bigl\{\, t \in [2025, 2075]
    \;\mid\; \text{elev}(x) - \text{MHHW} < \text{SLR}(t)\,\bigr\}.
  \label{eq:Tc}
\end{equation}
Parcels for which condition~\eqref{eq:chronic} is never satisfied
within the study period are assigned a sentinel value $T_c = 2076$
and labeled ``safe past 2075''.
The computation of $T_c$ was applied independently under the Low,
Intermediate, and High NOAA scenarios.

\subsection{Risk Scorer}

A logistic regression classifier was trained to predict whether a
parcel has high near-term flood risk, defined as $T_c^{\text{int}} \le 2050$
(i.e., the parcel reaches chronic inundation under the Intermediate
scenario before mid-century).

\textbf{Feature selection.}
Six physics-derived features were selected for the model:
(1)~\texttt{elev\_ft\_navd88} — ground-surface elevation,
(2)~\texttt{bfe\_ft} — FEMA base flood elevation,
(3)~\texttt{freeboard\_ft} — elevation minus BFE (negative means already
below the 100-year flood level),
(4)~\texttt{d\_100yr\_2025} — current 100-year flood inundation depth,
(5)~\texttt{d\_100yr\_2050\_int} — projected 100-year depth at 2050
under Intermediate SLR, and
(6)~\texttt{fld\_zone\_risk} — integer encoding of FEMA flood zone
(VE=4, AE=3, AH/AO=2, X~shaded=1, X~unshaded=0).
A seventh feature, \texttt{dist\_to\_coast\_m}, was computed as the
Euclidean distance from each parcel centroid to an approximated Biscayne
Bay shoreline LineString in EPSG:6437 metric coordinates.

Critically, the $T_c$ columns and the
\texttt{chronic\_by\_2050\_intermediate} flag were \textit{excluded}
from the feature matrix.
Including them would constitute data leakage, as the binary label is a
direct function of $T_c^{\text{int}}$.
The model therefore learns to approximate the chronic inundation
threshold using only observable physical characteristics, which is the
operationally useful form: a firm assessing a building for the first
time does not have pre-computed $T_c$ values.

\textbf{Model formulation.}
Features were standardized using \texttt{sklearn.StandardScaler}
($\mu = 0$, $\sigma = 1$).
Logistic regression with L2 regularization ($C = 1.0$, solver = \textsc{lbfgs},
max\_iter = 2000) was fitted on the full dataset after scaling.
The predicted probability of high near-term risk was scaled to a
0--100 risk score:
\begin{equation}
  \text{risk\_score}(x) \;=\; 100 \cdot \hat{p}(y=1 \mid \mathbf{x}),
\end{equation}
where $\hat{p}$ is the logistic sigmoid applied to the linear prediction.

\textbf{Evaluation.}
Model performance was estimated using 5-fold stratified cross-validation
(\texttt{StratifiedKFold}, random\_state~=~42) to preserve class balance
across folds.
Accuracy, ROC-AUC, and F1 score were computed on each held-out fold;
means and standard deviations are reported.
Because the purpose of the scorer is aggregation and prioritization
rather than predictive discovery, the final model was retrained on
all available data to maximize coefficient stability.


% ============================================================
\section{Results}
\label{sec:results}
% ============================================================

\subsection{SLR Scenarios}

Figure~\ref{fig:slr} shows the three NOAA 2022 SLR scenario curves for
Virginia Key linearly interpolated to annual resolution over 2025--2075.
By 2050, the Low, Intermediate, and High scenarios project cumulative
rises of \placeholder{0.4}, \placeholder{1.0}, and \placeholder{1.6}~ft
above the 2000 baseline, respectively.
The spread between Low and High at 2050 is
\placeholder{$\Delta = 1.2$}~ft—
a difference large enough to shift an entire generation of parcels
from the ``safe past 2075'' category under the Low scenario into the
``chronic by 2040'' category under the High scenario.

\begin{figure}[h]
  \centering
  \includegraphics[width=\columnwidth]{slr_scenarios.png}
  \caption{NOAA 2022 Sea-Level Rise scenarios for Virginia Key Station
    8723214 (Biscayne Bay, Miami FL), interpolated to annual resolution.
    Dashed vertical line marks 2050; dashed horizontal line marks the
    2000 baseline. Values at 2050 are annotated on each curve.}
  \label{fig:slr}
\end{figure}

\subsection{Inundation Analysis}

Figure~\ref{fig:tc} presents histograms of $T_c$ across all study
parcels for the three SLR scenarios.
Under the Intermediate scenario,
\placeholder{X}\% of parcels are projected to experience chronic
inundation before 2050, compared with \placeholder{X}\% under the
High scenario and \placeholder{X}\% under the Low scenario.
The histogram shifts markedly leftward from Low to High: parcels that
are ``safe past 2075'' under the Low scenario typically have
$T_c \in [2055, 2075]$ under High, reflecting how the accelerated
SLR trajectory advances the onset timeline by one to two decades for
parcels with modest elevation headroom above MHHW.

\begin{figure}[h]
  \centering
  \includegraphics[width=\columnwidth]{Tc_histograms.png}
  \caption{Distribution of chronic inundation onset year $T_c$ across
    study-area parcels under Low (left), Intermediate (center), and
    High (right) NOAA 2022 SLR scenarios.
    The sentinel value 2076 represents parcels that do not reach
    chronic inundation within the 2025--2075 study window.
    Vertical dashed lines mark 2050.}
  \label{fig:tc}
\end{figure}

\subsection{Risk Scorer Performance}

Five-fold stratified cross-validation yielded the following mean
performance metrics:
accuracy = \placeholder{0.XX} $\pm$ \placeholder{0.0X},
ROC-AUC  = \placeholder{0.XX} $\pm$ \placeholder{0.0X},
F1 score = \placeholder{0.XX} $\pm$ \placeholder{0.0X}.
These results indicate that the physics-derived features are highly
predictive of the binary $T_c^{\text{int}} \le 2050$ label, which is
expected given that the features are algebraically related to the
physics that produced the label; the cross-validation confirms that
the logistic approximation is stable across data splits.

Figure~\ref{fig:importance} shows the standardized logistic regression
coefficients sorted by absolute magnitude.
The two largest-magnitude coefficients are for
\placeholder{\texttt{elev\_ft\_navd88}} (negative, meaning higher ground
elevation strongly reduces predicted risk) and
\placeholder{\texttt{d\_100yr\_2050\_int}} (positive, meaning greater
projected 2050 inundation depth strongly increases predicted risk).
These findings are physically interpretable: elevation directly controls
the MHHW headroom in Eq.~\eqref{eq:Tc}, and projected 2050 flood depth
quantifies how much SLR has already ``consumed'' that headroom at
mid-century.

\begin{figure}[h]
  \centering
  \includegraphics[width=\columnwidth]{feature_importance.png}
  \caption{Standardized logistic regression coefficients (horizontal bars)
    for the seven flood-risk features.
    Positive coefficients (red) increase predicted risk; negative
    coefficients (blue) decrease it.
    Bars are sorted by absolute magnitude, placing the most influential
    features at the top.}
  \label{fig:importance}
\end{figure}

\subsection{Risk Map}

Figure~\ref{fig:map} presents the final parcel-level flood risk score
map for the Brickell study area.
The highest-risk parcels (score $> 67$, red) are concentrated along
the immediate waterfront south of the Miami River mouth, where a
combination of low natural ground elevation (typically $< 1.5$~ft
NAVD88), VE or AE flood zone classification, and minimal freeboard
above the FEMA BFE creates the most adverse combination of all six
model features.
Moving inland by as little as 150~m, risk scores drop to the
Medium range (33--67, orange--yellow) as ground elevations increase
and distance to the shoreline attenuates the contribution of
\texttt{dist\_to\_coast\_m}.
Parcels with Low risk scores ($< 33$, green) are predominantly those
on slightly elevated natural ridges or with post-2000 construction
that required higher finished-floor elevations under the 2014 and
2021 FEMA FIRM updates.

\begin{figure}[h]
  \centering
  \includegraphics[width=\columnwidth]{risk_map_static.png}
  \caption{Parcel-level flood risk scores (0 = low, 100 = high) for the
    Brickell waterfront study area, NOAA Intermediate SLR scenario.
    Background: CartoDB Positron; parcels colored by \texttt{RdYlGn\_r}
    continuous colormap. Scale bar represents 200~m; north arrow in
    upper right.}
  \label{fig:map}
\end{figure}


% ============================================================
\section{Discussion}
\label{sec:discussion}
% ============================================================

\subsection*{Engineering Relevance}

The risk score produced by this pipeline translates a complex
multi-scenario physics model into a single, auditable number that
a structural engineering firm can attach to a parcel record.
For a condominium board commissioning a 30-year reserve study—
now mandatory under Florida SB~4-D—a risk score above 67 signals
that the building's common-area waterproofing, drainage infrastructure,
and foundation inspection intervals should be planned against the
Intermediate or High SLR trajectory rather than the Low one.
This is not a substitute for site-specific geotechnical assessment,
but it enables a firm like M2E Consulting Engineers to triage a
portfolio of buildings—prioritizing milestone inspections for
properties whose chronic inundation onset is projected within a
single loan cycle ($\sim$10 years).

\subsection*{Limitations}

Five limitations constrain the current analysis.
(1)~\textit{DEM resolution}: The 1/3~arc-second 3DEP DEM has a nominal
horizontal resolution of $\sim$10~m and a vertical accuracy of
$\pm$0.3~m (1~$\sigma$) for open terrain, which exceeds the
$<$0.1~ft headroom differences that separate ``chronically flooded''
from ``safe'' parcels in the lowest-elevation decile of the study area.
(2)~\textit{Stillwater model}: Equation~\eqref{eq:depth} ignores wave
runup, surge, and groundwater infiltration, all of which contribute
to chronic inundation in practice.
(3)~\textit{Proxy shoreline}: The Biscayne Bay shoreline was approximated
as a single straight LineString at longitude $-80.187$, omitting
the bay's actual curvilinear geometry and the shielding effect of
Brickell Key.
(4)~\textit{Static BFE}: FEMA BFE values have not been updated for the
study area since the 2014 FIRM effective date and do not reflect
the additional surge from intermediate SLR.
(5)~\textit{Label circularity}: The binary label and several features are
algebraically related (they are both derived from the same BFE,
elevation, and SLR inputs), so cross-validation accuracy may
overestimate generalization performance on parcels outside the
study area.

\subsection*{Lessons Learned}

\textit{\textbf{[STUDENT TO COMPLETE]}} — This paragraph must be written
by the student based on their actual experience running the pipeline.
Describe: (a)~one aspect of the implementation that worked better than
expected, including which script and what the outcome was; (b)~one
specific failure or error encountered (e.g., a CRS mismatch, missing
file, API timeout, or unexpected data encoding), how it was diagnosed,
and how it was resolved; and (c)~one methodological choice that, in
hindsight, you would make differently—explaining both the original
choice and the preferred alternative.
This reflection should be 150--200 words and written in first person.


% ============================================================
\section{Conclusion}
\label{sec:conclusion}
% ============================================================

This paper presented a reproducible, open-source pipeline that couples
FEMA flood hazard data, USGS elevation, and NOAA sea-level rise
scenarios to a stillwater physics model and logistic regression risk
scorer, producing parcel-level chronic inundation onset timelines
and 0--100 flood risk scores for residential condominiums in
Brickell, Miami.
The scientific contribution is the demonstration that parcel-scale
$T_c$ estimates—previously confined to population-scale raster
analyses—can be derived from entirely public data at a spatial
resolution directly actionable for structural engineering decision
support under Florida SB~4-D.
Future work should replace the stillwater approximation with a
depth-averaged hydrodynamic model (e.g., ADCIRC or SLOSH) and extend
the spatial domain to all 330,000 condominium units in Miami-Dade
County.


% ============================================================
\begin{thebibliography}{9}
% ============================================================

\bibitem{sweet2022}
W.~V.~Sweet, B.~D.~Hamlington, R.~E.~Kopp, C.~P.~Weaver,
P.~L.~Barnard, D.~Bekaert, W.~Brooks, M.~Craghan,
G.~Dusek, T.~Frederikse, G.~Garner, A.~S.~Genz,
J.~P.~Gilford, E.~Hameeduddin, R.~Hetland,
C.~Johnson, R.~S.~Kadlec, J.~Larson, D.~Marcy,
J.~Marra, C.~Obeysekera, M.~Osler, E.~Pendleton,
D.~Roman, L.~Schmied, W.~Sweet, K.~D.~White,
and C.~Zuzak,
``2022 Sea Level Rise Technical Report,''
National Oceanic and Atmospheric Administration,
Silver Spring, MD, 2022.
[Online]. Available:
\url{https://oceanservice.noaa.gov/hazards/sealevelrise/sealevelrise-tech-report.html}

\bibitem{wdowinski2016}
S.~Wdowinski, R.~Bray, B.~P.~Kirtman, and Z.~Wu,
``Increasing flooding hazard in coastal communities
due to rising sea level: Case study of Miami Beach,
Florida,''
\textit{Ocean \& Coastal Management}, vol.~126,
pp.~1--8, 2016.

\bibitem{moftakhari2017}
H.~R.~Moftakhari, G.~Salvadori, A.~AghaKouchak,
B.~F.~Sanders, and R.~A.~Matthew,
``Compounding effects of sea level rise and fluvial
flooding,''
\textit{Proceedings of the National Academy of Sciences},
vol.~114, no.~37, pp.~9785--9790, 2017.

\bibitem{hauer2016}
M.~E.~Hauer, J.~M.~Evans, and D.~R.~Mishra,
``Millions projected to be at risk from sea-level rise
in the continental United States,''
\textit{Nature Climate Change}, vol.~6, no.~7,
pp.~691--695, 2016.

\bibitem{floridaSB4D}
Florida Legislature,
``Senate Bill 4-D: Building Safety,''
Florida Statutes \S{}718.112(2)(p), 2022.
[Online]. Available:
\url{https://www.flsenate.gov/Session/Bill/2022D/4D}

\bibitem{noaa_datums}
National Oceanic and Atmospheric Administration,
``Tidal Datums --- Virginia Key Station 8723214,''
CO-OPS Tidal Datum Table, 2023.
[Online]. Available:
\url{https://tidesandcurrents.noaa.gov/datums.html?id=8723214}

\bibitem{fema_nfhl}
Federal Emergency Management Agency,
``National Flood Hazard Layer (NFHL) Web Map Service,''
FEMA Flood Map Service Center, 2023.
[Online]. Available:
\url{https://msc.fema.gov/portal/home}

\end{thebibliography}

\end{document}
